\documentclass[11pt,a4paper]{article}

\usepackage{kotex}
\usepackage{fontspec}
\setmainfont{Times New Roman}
\setmainhangulfont{AppleMyungjo}
\setsanshangulfont{Apple SD Gothic Neo}

\usepackage[a4paper,top=18mm,bottom=18mm,left=20mm,right=20mm]{geometry}
\usepackage{fancyhdr}
\setlength{\headheight}{24pt}
\usepackage{enumitem}
\usepackage{mdframed}
\usepackage{array}
\usepackage{booktabs}

\setlength{\parindent}{1em}
\linespread{1.18}

\pagestyle{fancy}
\fancyhf{}
\renewcommand{\headrulewidth}{0.6pt}
\fancyhead[L]{\sffamily\bfseries 2026 고1 영어 변형 — 정답 및 해설 25}
\fancyhead[R]{\sffamily [자율주행차 신뢰도 도표]}
\renewcommand{\footrulewidth}{0pt}
\fancyfoot[C]{\framebox[16mm][c]{\thepage}}

\newmdenv[linewidth=0.6pt,roundcorner=3pt,innertopmargin=7pt,innerbottommargin=7pt,
          innerleftmargin=9pt,innerrightmargin=9pt,skipabove=5pt,skipbelow=5pt]{pbox}

\begin{document}

\begin{center}
{\fontsize{15}{17}\selectfont\sffamily\bfseries [지문 25] 정답 및 해설}
\end{center}
\vspace{8pt}

% 정답 표
\begin{center}
\begin{tabular}{|c|c|c|c|}
\hline
\textbf{문항} & \textbf{유형} & \textbf{정답} & \textbf{배점} \\\hline
1 & 도표 불일치 & ④ & 2점 \\\hline
2-(1) & 서술형 수치 단답 & 55\% & 3점 \\\hline
2-(2) & 서술형 수치 단답 & China & 3점 \\\hline
3 & 수치 추론 & India & 2점 \\\hline
4 & 서술형 수치 단답 & 48 percentage points & 2점 \\\hline
\end{tabular}
\end{center}

\vspace{12pt}

% Q1 해설
\noindent{\sffamily\bfseries 1번 해설} \quad \textbf{정답: ④}

\vspace{4pt}
\noindent\textbf{정답 근거:} 선지 ④는 ``All four countries showed the same percentage in the `a little' trust category''라고 주장한다.
그러나 표에 따르면 ``A little'' 비율은 U.K.\ 34\%, U.S.\ 30\%, China 35\%, India 35\%로, U.S.가 30\%로 나머지 나라와 다르다.
따라서 네 나라 모두 같다는 것은 사실이 아니므로 ④가 불일치 선지이다.

\vspace{6pt}
\noindent\textbf{오답 소거:}
\begin{itemize}[leftmargin=2em,itemsep=2pt,topsep=2pt]
  \item ① India만 ``A great deal''이 22\%로 20\% 이상이며, 나머지 국가(U.K.\ 4\%, U.S.\ 6\%, China 8\%)는 모두 한 자릿수. \textbf{일치} (정답 아님)
  \item ② U.K.\ 52\%, U.S.\ 51\% 모두 과반이 ``Not at all''. \textbf{일치} (정답 아님)
  \item ③ ``A lot'' 비율이 40\% 초과인 나라는 China(47\%)뿐. \textbf{일치} (정답 아님)
  \item ⑤ ``Not at all'' 비율: U.K.\ 52\%, U.S.\ 51\%, China 10\%, India 14\%. China가 10\%로 최저. \textbf{일치} (정답 아님)
\end{itemize}

\vspace{12pt}

% Q2 해설
\noindent{\sffamily\bfseries 2번 해설}

\vspace{4pt}
\noindent\textbf{(1) 모범답안:} \textbf{55\%}

\noindent\textbf{풀이:} China의 ``A great deal''(8\%) + ``A lot''(47\%) = \textbf{55\%}

\vspace{4pt}
\noindent\textbf{채점 포인트:} 55\% (단위 \% 포함 시 정답 인정; 계산 과정 불요)

\vspace{10pt}

\noindent\textbf{(2) 모범답안:} \textbf{China}

\noindent\textbf{풀이:} 각 나라의 ``A great deal'' + ``A lot'' 합산:
\begin{itemize}[leftmargin=2em,itemsep=1pt,topsep=2pt]
  \item U.K.: 4\% + 10\% = 14\%
  \item U.S.: 6\% + 13\% = 19\%
  \item China: 8\% + 47\% = 55\%
  \item India: 22\% + 29\% = 51\%
\end{itemize}
가장 높은 나라는 China(55\%)이다.

\noindent\textbf{채점 포인트:} ``China'' 또는 ``중국'' 정답 인정

\vspace{10pt}
\noindent{\sffamily\bfseries 3번 해설} \quad \textbf{정답: ④ India}
\par ``A great deal''과 ``A lot''의 합은 U.K. 14\%, U.S. 19\%, China 55\%, India 51\%이다. 가장 높은 나라는 China이고, 두 번째로 높은 나라는 India이다.

\vspace{8pt}
\noindent{\sffamily\bfseries 4번 해설} \quad \textbf{모범답안: 48 percentage points}
\par U.K.의 ``Not at all''은 52\%, ``A great deal''은 4\%이므로 차이는 52 - 4 = 48퍼센트포인트이다.

\end{document}
